Listening with a producer\textquotesingle s ears is a different activity from listening for pleasure. For
each of these, there is a specific thing to listen for. Put the track on, ignore everything
else, and follow just that one element for three minutes.

Almost all of this is on streaming. Where a record is hard to find, it is usually on
YouTube.

\subsection{The founding six}\label{li01_canon-the-founding-six}

Start here. Six records, roughly two hours including the long versions.

\textbf{Jesse Saunders, "On and On" (1984).} Usually cited as the first house record pressed to
vinyl. \emph{Listen for:} how little is happening. A drum machine, a bassline, a few stabs. This
is the entire genre in its most reduced form, and it will sound almost comically primitive.
That is the point. You can make this today in an afternoon.

\textbf{Marshall Jefferson, "Move Your Body" (1986).} \emph{Listen for:} the piano. Jefferson bought
a piano on credit and taught himself the riff by trial and error, and putting a real acoustic
instrument into a drum-machine record created piano house on the spot. Notice the riff is
simple enough that you could learn it in ten minutes, and notice that it does not matter.

\textbf{Mr. Fingers, "Can You Feel It" (1986).} Larry Heard made this on a Roland Juno-60 and a
TR-909 and almost nothing else. \emph{Listen for:} the space. Count how few elements there are at
any moment. The bassline is doing the emotional work and the pad is doing the atmosphere,
and there is no melody at all. This is the deep house template and it has never been
improved on.

\textbf{Phuture, "Acid Tracks" (1987).} \emph{Listen for:} the 303 line changing while never changing.
The notes are the same throughout; what moves is the filter cutoff and resonance. Twelve
minutes of one bassline, and it holds. This is the clearest possible demonstration that
timbre can carry a record when harmony and melody do not.

\textbf{Steve "Silk" Hurley, "Jack Your Body" (1986).} The one that went to number one in the UK.
\emph{Listen for:} the arrangement, specifically how elements drop in and out. It is a lesson in
how much mileage you get from muting things.

\textbf{Farley "Jackmaster" Funk, "Love Can\textquotesingle t Turn Around" (1986).} \emph{Listen for:} Darryl Pandy\textquotesingle s
vocal, which is enormous, gospel-derived, and completely uninterested in restraint. This is
the other pole of Chicago house: not machine music, but soul music that happens to have a
drum machine under it.

\subsection{The deep strand}\label{li01_canon-the-deep-strand}

Where the chords are. If you want to learn house harmony, steal from these.

\textbf{Mr. Fingers, "Mystery of Love" (1985).} \emph{Listen for:} the chord voicings. Seventh and
ninth chords, spaced widely, moving slowly.

\textbf{Fingers Inc., "Bring Down the Walls" (1986)}, with Robert Owens singing. \emph{Listen for:} how
the vocal and the track are equal partners rather than one sitting on top of the other.

\textbf{Ron Trent, "Altered States" (1990).} Trent was still a teenager. \emph{Listen for:} the length
and the patience. Nothing is in a hurry, and the track trusts you to stay with it.

\textbf{Larry Heard presents Mr. White, "The Sun Can\textquotesingle t Compare" (2005).} Proof that the man was
still doing it twenty years later. \emph{Listen for:} the bassline, which is a proper melodic
line rather than root notes.

\subsection{The anthems}\label{li01_canon-the-anthems}

House with something to say, usually about a better world. The gospel inheritance is
undisguised.

\textbf{Joe Smooth, "Promised Land" (1987).} \emph{Listen for:} the chord progression, which is
straightforwardly uplifting in a way most house avoids.

\textbf{CeCe Rogers, "Someday" (1987).} \emph{Listen for:} the piano and the vocal arrangement.

\textbf{Ten City, "Devotion" (1989).} \emph{Listen for:} the strings and the production polish. By
1989 house had proper budgets.

\textbf{Jamie Principle and Frankie Knuckles, "Your Love" (1986).} \emph{Listen for:} the arpeggio,
which is one of the most recognisable four bars in dance music, and how slowly everything
else moves underneath it.

\subsection{Acid and the harder side}\label{li01_canon-acid-and-the-harder-side}

\textbf{Adonis, "No Way Back" (1986).} \emph{Listen for:} the bassline and the sense of menace. Very
few elements, very dark.

\textbf{Armando, "Land of Confusion" (1988).} \emph{Listen for:} the relentlessness. This is where
acid house starts turning into techno.

\textbf{Lil Louis, "French Kiss" (1989).} \emph{Listen for:} the tempo breakdown in the middle, where
the whole record slows almost to a stop and then rebuilds. Nobody had done that. It works
because a dance floor will follow you anywhere if you commit.

\subsection{The nineties, when Chicago went global}\label{li01_canon-the-nineties-when-chicago-went-global}

\textbf{Cajmere, "Percolator" (1992).} \emph{Listen for:} the single sound the whole record is built
around. One weird bubbling noise, and that is the hook.

\textbf{Green Velvet, "Flash" (1995) and "La La Land" (2001).} Same person as Cajmere. \emph{Listen
for:} the spoken vocal used as a rhythmic element, and the deliberate weirdness. Curtis Jones
has a chemical engineering degree, which explains nothing but is a good fact.

\textbf{Paul Johnson, "Get Get Down" (1999).} \emph{Listen for:} the filtered disco loop. This is how
much of house is made: find four bars of an old record, filter it, put drums under it.

\textbf{DJ Sneak, "You Can\textquotesingle t Hide From Your Bud" (1996).} Same technique, more aggressive.

\subsection{Ghetto house and footwork}\label{li01_canon-ghetto-house-and-footwork}

The line from Chicago house to the present, and the part most outsiders never hear.

\textbf{DJ Funk} and \textbf{DJ Deeon}, anything on Dance Mania. \emph{Listen for:} how stripped it is.
Often just a 909, a bassline and a chant. No chords. It is house reduced to its skeleton and
it is much harder to make well than it sounds.

\textbf{RP Boo, "Baby Come On".} \emph{Listen for:} the triplet kick patterns and the enormous amount
of empty space. Footwork removes far more than it adds.

\textbf{DJ Rashad, \emph{Double Cup} (2013).} The album that took footwork worldwide. \emph{Listen for:}
how the vocal samples are chopped into fragments too short to carry meaning, and used purely
as rhythm. Rashad died in 2014.

\textbf{Traxman}, for the direct line back to the Trax catalogue.

\subsection{How to listen productively}\label{li01_canon-how-to-listen-productively}

Pick one track. Play it four times, and each time follow only one thing:

\begin{enumerate}
\def\labelenumi{\arabic{enumi}.}
\tightlist
\item
  \textbf{The kick and the drums.} Where does each element enter? What is the hi-hat doing?
\item
  \textbf{The bassline.} What notes, what rhythm, does it move or repeat?
\item
  \textbf{The arrangement.} Count bars. Write down what happens at bar 16, 32, 64. You will find
  almost everything lands on a multiple of 8.
\item
  \textbf{The space.} What is \emph{not} there? What did they leave out?
\end{enumerate}

\begin{listenbox}
\textbf{Listen for.} The fourth pass is the one that changes how you produce. Beginners add;
these records subtract. Ask what was left out, not what was put in.
\end{listenbox}

\subsection{A note on the sound quality}\label{li01_canon-a-note-on-the-sound-quality}

Much of this was recorded on cheap equipment, mixed in a bedroom, and pressed on recycled
vinyl by a label cutting corners. Some of it is genuinely badly made by any technical
standard.

It does not matter, and noticing that it does not matter is useful. Your first tracks will
also be technically poor. The question is whether the groove is right.
